\hypertarget{marked-correlation}{%
\section{Cancellation in the marked
correlation}\label{marked-correlation}}

We prove the estimate (B18) of
\protect\hyperlink{bilinear}{Section~\ref{bilinear}}. All the sums and
operators below are finite. We use prime Siegel--Walfisz, the
complete-cell discrepancy (C3), the high-frequency prime estimate
\protect\hyperlink{local-energy-the-high-frequency-prime-input}{(U1)},
smooth Fourier inversion, and the displayed finite energy estimates. The
bilinear exponential-sum bound comes from the row sums of an additive
Gram matrix and Cauchy--Schwarz.

Write \(L=\log x\) and \(\mathrm e(t)=e^{2\pi it}\). We retain the
actual marked groups and the multiplicity parity from
\protect\hyperlink{proof}{Section~\ref{proof}}. The BIG groups have
\(\log p\ge L^{3/10}\); the SMALL groups used for omissions have
\(\log p\le L^{1/5}\). Their harmonic masses lie in \([v,v_{\max}]\),
where \(v=\tfrac14\log(12/11)\) and \(v_{\max}=4\). We write these
masses as \(V_g\) in this chapter; they are \(M_g\) in the definition of
the weight. The fixed damping parameter is \(q=1/2\). The number of
groups is \(O(\log L)\). Constants implicit in this chapter may depend
on the fixed accuracy, frequency bounds, and mark counts, but not on
\(L\), the shift, or the endpoints of a fine cell.

The proof estimates a signed difference between the original correlation
and an arithmetic resampling term. A separate use of the fine-cell
discrepancy then makes the resampling term small. There are three
different spaces in this argument: a product of prime alphabets in
\protect\hyperlink{marked-correlation-orthogonal-averages-and-the-alternating-identity}{the
orthogonal-average argument},
\protect\hyperlink{marked-correlation-the-arithmetic-state-space-and-the-local-contraction}{the
residue and memory space}, and
\protect\hyperlink{marked-correlation-from-residue-averages-to-the-integer-interval}{the
integer-position space}. The local contraction is on the second space.
Its relation to the third is a moment identity and an interval-counting
estimate, not an identification of their norms. Throughout the operator
argument write \(s=2k\) for the integer shift in (B17). A forward frame
moves by \(sD\), and a reverse frame by \(-sD\).

\Needspace{6\baselineskip}

\hypertarget{marked-correlation-orthogonal-averages-and-the-alternating-identity}{%
\subsection{Orthogonal averages and the alternating
identity}\label{marked-correlation-orthogonal-averages-and-the-alternating-identity}}

In a BIG group \(P_g\), use the probability law

\[
V_g=\sum_{p\in P_g}\frac1p,
\qquad \mu_g(p)=\frac1{V_gp}.
\]

There are \(M=J+\ell\) labelled coordinates in each group. Here \(J\)
will be fixed later and the application has \(\ell=1\). Precisely, put
\(\mathcal X=\prod_{g\in\mathrm{BIG}}P_g^M\), with probability
\(\mu(x)=\prod_{g,i}\mu_g(x_{g,i})\), and use
\(\langle f,h\rangle_\mu=\sum_x\mu(x)\overline{f(x)}h(x)\). The tests
before symmetrization are arbitrary functions
\(F,H:\mathcal X\to\mathbb C\). Set

\[
(\mathcal SF)(x)=\frac1{(M!)^{G}}
 \sum_{\sigma\in\prod_g S_M}F((x_{g,\sigma_g(i)})_{g,i}),
\qquad f=\mathcal SF,\quad h=\mathcal SH,
\tag{M0}
\]

where \(G=\#\mathrm{BIG}\). Relabelling coordinates preserves \(\mu\),
so Jensen gives \(\|\mathcal SF\|_\mu\le\|F\|_\mu\). The actual kernel
independently averages its source and target permutations: changing
variables by their inverses moves these two averages onto \(F\) and
\(H\), exactly. This identity, rather than a symmetry assumption on the
arithmetic coefficients, licenses (M0). The tests can depend jointly on
all groups. The maps (T14)--(T15) identify them as Fourier coefficients
of the two physical endpoint tests.

Let \(E_i\) average coordinate \(i\). These are commuting orthogonal
projections. The Hoeffding components

\[
H_Sf=\prod_{i\in S}(I-E_i)\prod_{i\notin S}E_i f
\quad(S\subseteq[M])
\]

are orthogonal and sum to \(f\). Symmetry gives equal squared norms to
components of the same degree \(|S|\). Choose disjoint sets of \(m\)
probe coordinates and \(\ell\) trailing coordinates. At degree \(j\),
the fraction of components containing every probe and no trailing
coordinate is

\[
\frac{\binom{M-m-\ell}{j-m}}{\binom Mj}
\le\binom{m+\ell}{m}^{-1}.
\]

Consequently

\[
\left\|\prod_{i\in\mathrm{probe}}(I-E_i)
             \prod_{i\in\mathrm{tail}}E_i f\right\|_2^2
\le\binom{m+\ell}{m}^{-1}\|f\|_2^2. \tag{M1}
\]

Apply this successively to different groups, using Fubini. For
\(\ell=1\) and \(s\) groups the factor is \((m+1)^{-s}\). The actual BIG
groups contain two disjoint blocks \(\mathcal A,\mathcal B\), each of
size at least \(c\log L\), where \(c=(120\log2)^{-1}\). Given \(A_0>0\),
choose a fixed \(m\) so that \(A_0\le c\log(m+1)\). Then each block
gives squared-energy saving \(L^{-A_0}\).

The coupled phase cannot be moved through these coordinate averages. To
retain it exactly, put

\[
R(f,h)=\mathbb E[\overline f\,\Phi h],\qquad |\Phi|=1,
\]

For a fixed positive SMALL divisor \(d\), orientation \(o\in\{1,-1\}\),
and real \(\zeta,\theta\), its precise value is

\[
D(x)=d\prod_{g\in\mathrm{BIG}}\prod_{i=0}^{J-1}x_{g,i},
\qquad
\Phi(x)=D(x)^{io\zeta}\mathrm e(o sD(x)\theta).
\tag{M0a}
\]

Thus \(|\Phi|=1\), and \(\Phi\) ignores the \(\ell\) trailing
coordinates. It need not ignore the probes. Let \(U\) average all
trailing coordinates, and set

\[
E_a=U\prod_{i\in a}E_i,
\qquad P_{\mathcal A}=\sum_{a\subseteq\mathcal A\times[m]}
 (-1)^{|a|}E_a
 =U\prod_{i\in\mathcal A\times[m]}(I-E_i).
\]

Define \(P_{\mathcal B}\) in the same way. Expanding the two products
gives the exact identity

\[
\begin{aligned}
R(Uf,Uh)&-
\sum_{\substack{\varnothing\ne a\subseteq\mathcal A\times[m]\\
                 \varnothing\ne b\subseteq\mathcal B\times[m]}}
 (-1)^{|a|+|b|}R(E_af,E_bh)\\
&=R(P_{\mathcal A}f,Uh)+R(Uf,P_{\mathcal B}h)
  -R(P_{\mathcal A}f,P_{\mathcal B}h).
\end{aligned} \tag{M2}
\]

Both subsets in the comparison sum are nonempty. Each term on the right
has a contracted endpoint; since \(|\Phi|=1\), Cauchy--Schwarz and (M1)
give

\[
|\text{right side of (M2)}|^2
 \le9L^{-A_0}\|f\|_2^2\|h\|_2^2. \tag{M3}
\]

The number of comparison patterns is a fixed power of \(L\): \(m\) is
fixed and there are \(O(\log L)\) groups. We will pay this cost in the
arithmetic approximation to each pattern. It is not subtracted from the
exponent supplied by the already summed identity (M2).

\Needspace{6\baselineskip}

\hypertarget{marked-correlation-implementing-the-averages-by-characters}{%
\subsection{Implementing the averages by
characters}\label{marked-correlation-implementing-the-averages-by-characters}}

For a nonempty collection of probes, let \(D\) be their product under
the product probability law. Partition its logarithm into cells
\(I_j=[e^{jh},e^{(j+1)h})\), put \(\nu_j=\Pr(D\in I_j)\), and retain
cells with \(\nu_j\ge\xi\). The arithmetic kernel is

\[
K(D,X)=\sum_{\nu_j\ge\xi}
 \frac{\mathbf1_{D,X\in I_j}}{\nu_j}
 \sum_{\substack{\chi\ \mathrm{primitive}\\\operatorname{cond}\chi\le Q}}
       \overline{\chi(D)}\chi(X). \tag{M4}
\]

The reference operation is exact matching of the product, not the
constant kernel. If \(\pi(X)=\Pr(D=X)\), put

\[
I_\pi(D,X)=\frac{1_{D=X}}{\pi(X)},\qquad
\sum_D\pi(D)I_\pi(D,X)\phi(D)=\phi(X)
\quad(\pi(X)>0).
\tag{M4a}
\]

Set the kernel to zero when \(\pi(X)=0\); such \(X\) do not occur in the
product law. Applying (M4a) separately in the two probe blocks sets the
auxiliary products equal to the target first-probe and source
second-probe products. The other, freed endpoint coordinates are then
averaged with their original laws. This is the ordinary resampling form
in (M2). In particular, it evaluates the phase at the spliced products;
it does not average that phase independently. The estimate below
compares this matching form with (M4), including its phase and endpoint
tests; it does not assert that their kernels are close in operator norm.

All marked primes exceed \(Q\). Characters in (M4) thus have modulus one
on its support. The cell Gram matrix has diagonal \(\nu_j\). For two
different primitive characters, freeze every prime coordinate but one.
The condition on the product clips that remaining prime to an interval.
The quotient of the two characters induces a nonprincipal primitive
character with conductor at most \(Q^2\). Prime Siegel--Walfisz and
partial summation bound the corresponding entry by

\[
|G_j(\chi,\chi')|\le (K_1/V_g)L^{-B_1}.
\]

The other coordinates have total probability one. Summing over them does
not introduce the number of tuples. There are at most \(Q^2\)
characters. Therefore, if \(Q\le L^{C_q}\) and
\(\xi^{-1}\le L^{C_\xi}\), requesting accuracy \(B_1=B+C_\xi+2C_q\)
gives

\[
\|K\|_{2\to2}\le1+(K_1/v)L^{-B}. \tag{M5}
\]

Different cells give orthogonal blocks of this operator, so their number
does not multiply (M5).

To compare (M4) with ordinary resampling, retain the full phase
\(\mathrm e(\theta'XY)(XY)^{i\zeta}\). On a major arc
\(\theta'=a/r+\beta\), \(r\le Q\), the function \(\mathrm e(aXY/r)\) on
units is a sum of primitive-character modes with coefficient \(\ell^1\)
norm at most \(r\). The two kernels reproduce these modes up to

\[
(2\epsilon+\epsilon^2)r,
\qquad \epsilon=(K_1/v)L^{-B}.
\]

Within two matched logarithmic cells, the remaining phase varies by at
most

\[
Q^4\bigl(|\zeta|+2\pi|\beta|R_0\bigr)(h_A+h_B),
\qquad R_0=e^{(j_A+1)h_A+(j_B+1)h_B}. \tag{M6}
\]

Here \(R_0\) bounds the product of the two cell variables. This follows
by differentiating the phase in the two logarithmic coordinates. The
factor \(Q^4\) bounds the two finite character sums.

On minor arcs, move the two kernels onto the endpoint tests and use
(M5), eventually with norm at most \(2\). The finite integer bilinear
exponential-sum inequality on dyads of sizes \(M_0,N_0\) has spectral
factor bounded by

\[
\frac1r+\frac2{N_0}+\frac{\log(2r)}{M_0}
       +\frac{r\log(2r)}{M_0N_0}. \tag{M7}
\]

For example, for windows of integer lengths \(P\le M_0\), \(Q'\le N_0\),
the bound preceding (M7) is \([P+(1+Q'/r)(P+r\log(2r))]/(M_0N_0)\). Thus
(M7) is at most \(5L^{-D}\) provided

\[
r,N_0\ge L^D,\quad M_0\ge L^D\log(2r),\quad
M_0N_0\ge L^Dr\log(2r). \tag{M8}
\]

Here is a denominator choice that verifies this alternative. Put
\(d_*=2/5\), \(S_D=L^{D+d_*+2}\), and apply Dirichlet approximation with
cutoff \(\lfloor M_0N_0/S_D\rfloor\). It gives a reduced \(a/r\) with

\[
r\le M_0N_0/S_D,\qquad
|\beta|M_0N_0\le S_D/r,\qquad |\beta|\le r^{-2}.
\]

Each nonempty BIG probe product is at least \(e^{L^{3/10}}\) and has
logarithm \(O_m(L^{2/5}\log L)\); in particular
\(\log(2M_0N_0)\le L^{d_*+2}\) eventually. If \(r>Q\ge L^D\), these
bounds give all four inequalities in (M8). If \(r\le Q\), the matched
cells satisfy \(R_0\le4e^{h_A+h_B}M_0N_0\), so \(|\beta|R_0\) has a
fixed polynomial bound. This supplies the major-arc variation estimate
with explicit control of its denominator.

Passing from labelled prime tuples to their products costs at most

\[
\prod_g\frac{k_g!}{V_g^{k_g}}, \tag{M9}
\]

because a product in a group has at most \(k_g!\) ordered
representations. Here \(k_g\le m\), so (M9) is another fixed power of
\(L\). After including it in the requested minor-arc precision, the
source bounds for the raw term and the two-kernel term are,
respectively, \(125L^{-B}\|f\|_2^2\|h\|_2^2\) and
\(2000L^{-B}\|f\|_2^2\|h\|_2^2\). Their squared difference is bounded by
\(4250L^{-B}\|f\|_2^2\|h\|_2^2\).

If the two product ranges meet \(n_A,n_B\) cells, the probability lost
on discarded cells is at most \(2(n_A+n_B)\xi\). To bound their number,
each probe product has logarithm at most \(B_L=mC_G L^{2/5}\log L\),
where \(C_G\log L\) bounds the BIG group count. Consequently

\[
n_A,n_B\le2+B_L/h,\qquad
\#\{\text{coarse dyads for either product}\}
 \le2+B_L/\log2.
\tag{M9a}
\]

These are polynomial bounds for the mesh (M10), uniform in the
conditioned shared coordinates. They apply also when the product range
is further clipped. The number of patterns is at most
\(2^{m(|\mathcal A|+|\mathcal B|)}\), also polynomial. Shared
coordinates are conditioned on and then integrated by Cauchy--Schwarz.
For a requested one-rectangle saving \(B_{\rm save}\), let
\(C_{\rm factor}\) pay for (M9). The following choices suffice:

\[
\begin{gathered}
D\ge B_{\rm save}+2C_{\rm factor},\qquad
C_s\ge D+d_*+2,\\
C_4\ge C_s+1,\quad C_4\ge C_{\rm freq},\\
H_m\ge B_{\rm save}+4C_q+C_4,\qquad
C_\xi\ge B_{\rm save}+H_m+d_*+4.
\end{gathered}
\]

Here \(|\zeta|\le L^{C_{\rm freq}}\) and \(C_q\) bounds the conductor
cutoff. Request precision \(B_{\rm save}+C_q\) in the rational-phase
reproduction to pay for \(r\le L^{C_q}\); the prime Siegel--Walfisz
precision is then increased by \(C_\xi+2C_q\) as in (M5). Any fixed
constants can be absorbed by an extra power and a larger threshold.
After paying the pattern and dyad counts in \(B_{\rm save}\), (M3)--(M9)
give the small local signed-operator estimate with

\[
h=\frac{\log2}{\lceil4L^{H_m}\rceil},\qquad
Q=\lceil4L^{D+1}\rceil,\qquad \xi=L^{-C_\xi}. \tag{M10}
\]

Choose \(C_q>D+1\) after \(D\) and before \(H_m,C_\xi\).

\Needspace{6\baselineskip}

\hypertarget{marked-correlation-the-arithmetic-state-space-and-the-local-contraction}{%
\subsection{The arithmetic state space and the local
contraction}\label{marked-correlation-the-arithmetic-state-space-and-the-local-contraction}}

\Needspace{6\baselineskip}

\hypertarget{marked-correlation-residues-active-marks-and-memory}{%
\subsubsection{Residues, active marks and
memory}\label{marked-correlation-residues-active-marks-and-memory}}

Fix \(N=2\lceil L^{13/25}\rceil\), \(M=J+1\), and finite cutoffs \(B_g\)
for the BIG memory lengths. The estimates below are uniform in these
cutoffs. For a SMALL group put

\[
\mathbb T_g=\prod_{p\in P_g}\mathbb Z/p\mathbb Z,
\qquad
\mathscr S_g=\{(y,a):y\in\mathbb T_g,\ a:[M]\hookrightarrow P_g,
                         \ y_{a_i}=0\ (i\in[M])\}.
\tag{T1}
\]

Here and below \([M]=\{0,\ldots,M-1\}\). Give every element of
\(\mathscr S_g\) mass \(|\mathbb T_g|^{-1}V_g^{-M}\), and write
\(\omega_g(y)=\#\{p:y_p=0\}\). This is Haar probability in \(y\) and
counting measure in the valid ordered marks, with the displayed factor
\(V_g^{-M}\); it is not a probability on the marked fibre. Translation
means \(y_p\mapsto y_p+\Delta\pmod p\) simultaneously for every \(p\).

For BIG primes set

\[
q_i=\begin{cases}\sqrt q,&i=0,N,\\q,&0<i<N,\end{cases}
\quad \eta_i=1-q_i,\quad H_N=\sum_{i=0}^N\eta_i,
\quad b_p=1-H_N/p,\quad \widetilde\mu_g(p)=\frac1{V_gp b_p}.
\tag{T2}
\]

Eventually \(b_p\ge1/2\). A BIG state consists of an active array
\(a_g\in P_g^M\) and an ordered list \(((p_1,r_1),\ldots,(p_j,r_j))\),
\(0\le j\le B_g\), where \(r_i\in\mathbb Z/p_i\mathbb Z\). Its mass is

\[
\left(\prod_{i<M}\widetilde\mu_g(a_{g,i})\right)
 \frac1{j!}\prod_{i=1}^j\widetilde\mu_g(p_i).
\tag{T3}
\]

The residues in this list have \textbf{counting} measure: summing the
mass of \((p,r)\) over \(r\) gives \(p\widetilde\mu_g(p)\), while its
zero residue has mass \(\widetilde\mu_g(p)\). Equal primes and equal
atoms are allowed here. The active BIG array also has no distinctness
condition at this stage. Distinctness is imposed on the appropriate
birth lists in the transition and then globally in the forest expansion.

Let \(\widetilde{\mathscr S}\) be the product of (T1) and these BIG
spaces, with mass \(w\). Let \(\pi\) replace each ordered OLD list by
its multiset, leaving all SMALL states and active arrays unchanged. On
\(\mathscr S=\pi(\widetilde{\mathscr S})\) put

\[
\bar w(X)=\sum_{\pi(x)=X}w(x),\qquad
\mathcal E(F)=\sum_X\bar w(X)|F(X)|^2,
\qquad
(\Pi u)(X)=\frac{\sum_{\pi(x)=X}w(x)u(x)}{\bar w(X)}.
\tag{T4}
\]

A zero-mass fibre is assigned zero. In our nonempty fibres the masses
are positive. Pullback \(\pi^*F=F\circ\pi\) is an isometry,
\(\Pi\pi^*=I\), and Jensen gives
\(\mathcal E(\Pi u)\le\sum_xw(x)|u(x)|^2\). For a multiset with atom
multiplicities \(n_z\), the OLD part of its mass is
\(\prod_z\widetilde\mu_g(p(z))^{n_z}/\prod_z n_z!\). This follows by
counting its \(j!/\prod_z n_z!\) ordered representatives in (T3). In
particular, passing to multisets never discards collisions.

An OLD atom records a previously introduced prime and its current
residue. On a step \(\Delta=o sD\) all surviving residues become
\(r+\Delta\pmod p\). A \emph{promotion} supplies an unshared target
active slot from an OLD occurrence whose translated residue is zero;
distinct promoted slots use distinct OLD indices. A \emph{store} puts an
unshared source active occurrence into OLD, at residue
\(\Delta\pmod p\). Stored and promoted primes must not divide \(D\). The
remaining unshared target slots are new births sampled with
\(\widetilde\mu_g\). The output length must stay below \(B_g\). A
transition is \emph{clean} when it has no stores and no promotions.
Clean transitions translate every OLD occurrence, including repetitions.

For definiteness the active operator is the following finite sum. Choose
the raw pattern or a nonempty pair \((a,b)\) of probe subsets, auxiliary
primes with the original harmonic law \(\mu\), and two independent
uniform permutations in each SMALL and BIG group. Copy the shared active
coordinates, use the divisor \(D\) of (M18b), and sum all legal
store/promotion choices just described. In each BIG group the
coefficient is

\[
\rho_0^{z_g+z'_g}V_g^{-t_g}
 \prod_{i\text{ fresh at the target}}\widetilde\mu_g(a'_{g,i}),
\tag{T5}
\]

where \(z_g,z'_g\) count zero-residue OLD occurrences at the two
endpoints and \(t_g\) is the number of promotions. Choose
\(q^{1/4}<\rho_0<1\). In each SMALL group multiply by
\(V_g^{-1}q^{(\omega_g(y)+\omega_g(y')-2M)/2}\), imposing shared marks
and \(y'=y+\Delta\). Finally multiply by the auxiliary mass, the
permutation masses, and \(K_fD^{i\zeta}\mathrm e(sD\theta)\) in the
forward orientation; the reverse orientation conjugates
\(K_fD^{i\zeta}\) and changes \(\theta\) to \(-\theta\) in the return
phase. There is only one such complex scalar for the full frame.

The operator \(E_o\) is this row sum, pulled back and averaged by (T4),
with the two gates \(|\mathrm{OLD}|\le T=\lceil L^2\rceil\) and the
condition that the list of fresh target numerical primes is distinct.
Only distinctness between different birth times is omitted. These are
the operators denoted \(E_i\) in (F7), with alternating \(o\). The
integer-position matrix (M18c) has a different carrier and will be
related to their full path sum in
\protect\hyperlink{marked-correlation-from-residue-averages-to-the-integer-interval}{the
residue-to-interval transfer}.

\Needspace{6\baselineskip}

\hypertarget{marked-correlation-the-frozen-big-form-including-its-errors}{%
\subsubsection{The frozen BIG form, including its
errors}\label{marked-correlation-the-frozen-big-form-including-its-errors}}

We first identify the harmonic form bounded by (M1)--(M10). Freeze the
SMALL shared tuple and its product \(d>0\). For a pattern \(\nu\), let
\(S_g\subset[J]\) be its shared coordinates. For source and target
arrays \(u,v\) and permutations \(\sigma,\tau\), put

\[
B_\mu(u,v)=\prod_g\left[
 \prod_{i\in S_g}1_{u_{g,\sigma_g(i)}=v_{g,\tau_g(i)}}
 \prod_{i\notin S_g}\mu_g(v_{g,\tau_g(i)})\right].
\tag{T6}
\]

In this formula the complement is in \([M]\), so it includes the
trailing coordinate. Multiply (T6) by the auxiliary mass, the two
permutation masses, the two-endpoint auxiliary ban, and the single
oriented scalar from (T5); sum over \(\nu\), auxiliaries and
permutations to obtain a kernel \(\mathsf B_\mu(d,\zeta,\theta,o,s)\).
Its pairing is
\(\sum_{u,v}\mu(u)\overline{F(u)}\mathsf B_\mu(u,v)H(v)\). This defines
the frozen form for arbitrary tests; no factorization of \(F\) or \(H\)
between groups is required.

Here is the identity connecting it to (M2). Temporarily remove the
auxiliary ban. In a comparison pattern replace the two character kernels
by ordinary probe resampling. Use the shared-copy indicators to identify
the unchanged coordinates. At a probe belonging to the first subset use
the target value, at a probe belonging to the second use the source
value, and integrate the freed counterpart. The trailing coordinates at
both endpoints are independently integrated. The divisor in the phase is
the head product of this spliced array, so the resulting pairing is
exactly \(-(-1)^{|a|+|b|}R(E_a\mathcal SF,E_b\mathcal SH)\). The empty
pattern has the opposite sign and gives
\(R(U\mathcal SF,U\mathcal SH)\). Changing variables by the independent
endpoint permutations proves

\[
\mathsf B_\mu(F,H)=\mathsf R(F,H)
              +\mathsf C(F,H)-\mathsf Z(F,H),
\tag{T7}
\]

where \(\mathsf R\) is the right side of (M2), \(\mathsf C\) is the sum
of (character-resampling minus ordinary-resampling) comparison forms
without the ban, and \(\mathsf Z\) is the character form on auxiliaries
meeting either endpoint. This is an identity of signed forms, before
taking absolute values. In particular, the phase is never commuted
through a coordinate average.

Equations (M4)--(M10) bound \(|\mathsf C|^2\) by
\(L^{-A}\|F\|_\mu^2\|H\|_\mu^2\) at any fixed accuracy \(A\). The same
bound holds for \(\mathsf Z\) eventually: for a fixed endpoint array,
the probability that one of at most \(mG\) auxiliaries hits its \(MG\)
marked values is at most \(mMG^2/(v e^{L^{3/10}})\). Split the ban into
its source and target parts and apply weighted Cauchy--Schwarz to each
part. The unrestricted marginal at the other endpoint is the original
product law. The pointwise kernel cap and pattern sum are polynomial in
\(L\), so this exponential atom saving absorbs them. This argument keeps
the tests in their original \(L^2\) norms; it does not bound a test on a
rare set by its supremum. Together with (M3) and
\(|a+b-c|^2\le3(|a|^2+|b|^2+|c|^2)\), (T7) gives the explicit squared
bound \(33L^{-A}\) times the two energies. To obtain exponent \(2E\),
use \(A=2E+2\). The reduction to \(d=s=1\) replaces \(\theta\) by
\(ds\theta\) and introduces the unit scalar \(d^{io\zeta}\); hence this
bound is uniform in \(d>0\), \(s\ne0\) and \(\theta\in\mathbb R\).

The density in (T6) must next be changed to (T2). This is done with one
multiplier common to every pattern and permutation. Put

\[
r_g(p)=b_p^{-1/2},\qquad
(\mathcal AF)(u)=\prod_{g,i}r_g(u_{g,i})F(u).
\tag{T8}
\]

Then \(\|\mathcal AF\|_\mu=\|F\|_{\widetilde\mu}\). The precise ratio of
the tilted joint law to the harmonic joint law, after this change at
both endpoints, is

\[
R_\nu(u,v;\sigma,\tau)=
 \prod_g\prod_{i\notin S_g}
 r_g(u_{g,\sigma_g(i)})r_g(v_{g,\tau_g(i)}).
\tag{T9}
\]

For a shared coordinate the indicator in (T6) makes the two endpoint
square roots supply its full ratio; for a free coordinate an additional
square root is needed at each end. This verifies (T9) coordinate by
coordinate. Auxiliary draws still use \(\mu\), not \(\widetilde\mu\).

Let \(\mathsf B_{\widetilde\mu}^{\rm nd}\) include the fresh-birth list
condition in \(E_o\), and let \(\mathsf B_\mu^{\rm bad}\) be the same
harmonic form restricted to its failure. The exact identity is

\[
\mathsf B_{\widetilde\mu}^{\rm nd}(F,H)
 =\mathsf B_\mu(\mathcal AF,\mathcal AH)
   +\mathsf T(F,H)-\mathsf B_\mu^{\rm bad}(\mathcal AF,\mathcal AH),
\tag{T10}
\]

where \(\mathsf T\) uses the multiplier \((R_\nu-1)\) and retains the
fresh-birth condition. Set \(P=e^{L^{3/10}}\), \(a_0=1-H_N/P\),
\(I=MG\), and

\[
\delta_{\rm tilt}=a_0^{-I}-1,\qquad
\mathcal K=1+(2^{m|\mathcal A|}-1)(2^{m|\mathcal B|}-1)Q^4\xi^{-2}.
\tag{T11}
\]

The raw pattern has cap one and every other pattern has cap
\(Q^4\xi^{-2}\), so \(\mathcal K\) is their total cap. From (T9),
\(1\le R_\nu\le a_0^{-I}\); normalized auxiliary and permutation sums
and the two marginals of (T6) therefore give

\[
\begin{aligned}
|\mathsf T(F,H)|^2
 &\le(\mathcal K\delta_{\rm tilt})^2
       \|F\|_{\widetilde\mu}^2\|H\|_{\widetilde\mu}^2,\\
|\mathsf B_\mu^{\rm bad}(\mathcal AF,\mathcal AH)|^2
 &\le\mathcal K^2\frac{I^2}{vP}
       \|F\|_{\widetilde\mu}^2\|H\|_{\widetilde\mu}^2.
\end{aligned}\tag{T12}
\]

For the second bound, expose the fresh coordinates in (T6). Any pair
coincides with conditional probability at most \(1/(vP)\); there are at
most \(I^2\) pairs. The restricted row has this saving and the full
column has mass one. Weighted Schur gives the displayed squared bound,
valid even for tests concentrated on the bad set. Since \(H_N\le5L\),
\(I=O_J(\log L)\), and \(\mathcal K\le L^{C(m,C_q,C_\xi)}\), both errors
in (T12) are smaller than any inverse power of \(L\). Applying the
three-term square inequality to (T10) gives \(9L^{-2E}\) times the
tilted endpoint energies, after choosing the harmonic precision. This is
the full tilted BIG estimate, including the ban and the actual list
filter.

\Needspace{6\baselineskip}

\hypertarget{marked-correlation-the-small-maps-and-the-translated-old-residues}{%
\subsubsection{The SMALL maps and the translated OLD
residues}\label{marked-correlation-the-small-maps-and-the-translated-old-residues}}

It remains to transfer this bound to (T4). Fix the OLD lengths \(j_g\)
and the ordered OLD primes \(\boldsymbol\pi_g\), retaining their
residues as variables. On this sector use the finite abelian group

\[
\mathbb T=\prod_{g\in\mathrm{SMALL}}\mathbb T_g
       \times\prod_{g\in\mathrm{BIG}}\prod_{i=1}^{j_g}
                                  \mathbb Z/\pi_{g,i}\mathbb Z.
\tag{T13}
\]

Repeated OLD primes give separate coordinates. The translation
\(t\mapsto t+\Delta\) acts diagonally on every occurrence; this is
exactly the clean OLD update. For a SMALL permutation \(\sigma\) and
shared tuple \(d=(d_{g,i})_{g,i<J}\), define

\[
(P_{\sigma,d}U)(a;y,r)=
 \prod_{g\in\mathrm{SMALL}}
    V_g^{-1/2}q^{\max(\omega_g(y)-M,0)/2}
 \sum_{\substack{(y,c)\in\prod_g\mathscr S_g\\
                   c_{g,\sigma_g(i)}=d_{g,i}\ (i<J)}}
          U((y,c),a,r).
\tag{T14}
\]

A fibre with no valid marks contributes zero. These are conditional
sums, not conditional probabilities. In the clean coefficient the SMALL
scalar is exactly the product of the two half factors in (T14). The
shared tuple has the same value at the two endpoints, and their residues
differ by \(\Delta=o sD\). Partitioning both endpoint sums by
\((d,y,r)\) gives precisely the pairing of \(P_{\sigma,d}U\) with the
translated \(P_{\tau,d}F\), with the frozen BIG kernel between them. The
displacement still depends on \(d\), the BIG shared marks and the
auxiliaries. It has not been frozen before this disintegration.

Use normalized Fourier coefficients, so that

\[
\widehat u(\chi)=\frac1{|\mathbb T|}\sum_tu(t)\overline{\chi(t)},
\qquad
\widehat{u(\cdot+\Delta)}(\chi)=\chi(\Delta)\widehat u(\chi).
\]

Parseval replaces translation by multiplication by a character. Write
\(\chi(\Delta)=\mathrm e(\vartheta_\chi\Delta)\); the frozen BIG block
now has frequency \(\theta+\vartheta_\chi\). The bound above is uniform
in this frequency, however large the OLD sector. For this fixed prime
sector put

\[
c_{\boldsymbol\pi}=
 \prod_{g\in\mathrm{SMALL}}V_g^{-M}
 \prod_{g\in\mathrm{BIG}}
       \frac{\prod_{i=1}^{j_g}\pi_{g,i}\widetilde\mu_g(\pi_{g,i})}{j_g!}.
\]

The factors \(\pi_{g,i}\) convert OLD counting measure to normalized
Haar measure on (T13). The exact clean pairing in this sector is

\[
c_{\boldsymbol\pi}\,
 \mathbb E_{\sigma,\tau}\sum_d\sum_{\chi\in\widehat{\mathbb T}}
 \mathsf B_{\widetilde\mu}^{\rm nd}
   \bigl(\widehat{P_{\sigma,d}U}(\chi),
         \widehat{P_{\tau,d}F}(\chi);
         d_S,\zeta,\theta+\vartheta_\chi,o,s\bigr),
\quad d_S=\prod_{g,i<J}d_{g,i}.
\tag{T15}
\]

The two arguments of \(\mathsf B\) are functions of the whole BIG active
array \(a\). After (T8) and (M0) they are the tests \(f,h\) of the
orthogonal-average argument. Formula (T15) is the required operator
identity: it keeps the full joint phase, both endpoint permutations, all
shared values, and every translated OLD occurrence.

There is no loss equal to \(|\mathbb T|\) in its norm estimate. For a
fixed \(y,d\), the number of completions in one SMALL group is
\((\omega_g(y)-J)_{1}\); for \(\ell\) omitted marks it is
\((\omega_g(y)-J)_{\ell}\). Cauchy--Schwarz in that fibre, followed by
summation over \(d\) and Parseval, bounds each lifted endpoint energy by
\(C_{\rm cond}\) times its original sector energy, where

\[
C_{\rm cond}=\prod_{g\in\mathrm{SMALL}}
       \frac{\ell!}{(1-q)^{\ell+1}V_g^\ell},\qquad \ell=1.
\tag{T16}
\]

Indeed, on a nonempty fibre \(\omega_g=J+\ell+z\), and
\(q^z(z+\ell)_{\ell}\le\sum_{z\ge0}q^z(z+\ell)_{\ell} =\ell!/(1-q)^{\ell+1}\).
Every original ordered mark state occurs once when its shared tuple is
summed. Cauchy--Schwarz over \((d,\chi)\) and then over sectors yields

\[
|\langle U,E_{\rm clean}^{(1)}F\rangle|^2
 \le9L^{-2E_0}C_{\rm cond}^2\,\mathcal E(U)\mathcal E(F).
\tag{T17}
\]

Here the superscript means that the OLD endpoint parameter is one. For
general \(\rho_0\le1\), multiply both tests by
\(\rho_0^{\#\text{zero OLD}}\); these multipliers and both OLD gates are
contractions. Pullback and orbit averaging in (T4) are contractions as
well. To pass from a pairing bound to an operator bound, set the first
test equal to the output of the operator and cancel its norm (the zero
output is immediate).

If \(C_G\log L\) bounds the number of SMALL groups and \(V_g\ge v\),
then \(C_{\rm cond}\le L^{c_{\rm cond}}\), with
\(c_{\rm cond}=C_G\log\max(1,\ell!/((1-q)^{\ell+1}v^\ell))\). Choose
\(E_0=E+c_{\rm cond}+1\) before choosing \(m\) or \(J\). For \(L\ge3\),
(T17) is at most \(L^{-2E}\) times the two energies. This proves the
clean part of (F7), with all normalizations specified.

\Needspace{6\baselineskip}

\hypertarget{marked-correlation-promotions-stores-and-the-full-transition}{%
\subsubsection{Promotions, stores and the full
transition}\label{marked-correlation-promotions-stores-and-the-full-transition}}

Fix an OLD BIG prime \(p\), its residue \(r\), the BIG divisor \(D_B\),
and an orientation \(o\in\{\pm1\}\). Let \(F\) be the product of all
SMALL marks and \(U\) the product of the omitted \(\ell\) marks in each
SMALL group. The retained SMALL divisor is \(D_S=F/U\). Promotion
requires

\[
r+osD_SD_B\equiv0\pmod p,
\quad\text{or equivalently}\quad
rU+osD_BF\equiv0\pmod p. \tag{M11}
\]

The constant term is nonzero: \(p\nmid D_BF\) and \(0<|s|<p\).
Furthermore

\[
\log U\le\ell\,\#\mathrm{SMALL}\,L^{1/5}
 =O(\log L\,L^{1/5})<L^{3/10}\le\log p.
\]

Thus \(0<U<p\), and (M11) determines at most one \(U\). Distinct marks
within each group and disjointness of the groups determine the omission
tuple from \(U\). Uniform coordinate permutations induce uniform
omission subsets. Hence

\[
\Pr(\text{an OLD promotion})
\le\frac{\#\mathrm{OLD}}
 {\binom{J+\ell}{\ell}^{\#\mathrm{SMALL}}}. \tag{M12}
\]

In the application \(\ell=1\), \(\#\mathrm{OLD}\le\lceil L^2\rceil\),
and \(\#\mathrm{SMALL}\ge c_0\log L\) with \(c_0=(40\log2)^{-1}\). The
right side is at most \(2L^{2-c_0\log(J+1)}\). If \(L^\kappa\) bounds
the complete row cost, choose \(J\ge m\) so that

\[
2H_0+2\kappa+3\le c_0\log(J+1). \tag{M13}
\]

This makes the promotion row rare enough for the squared-energy saving
\(L^{-2H_0}\). Reverse the transition orientation to obtain the same
estimate for storing an active label: it is a rare column. The weighted
Schur inequality combines each rare row or column bound with the full
complementary row bound, then adds the two exceptional operators.
Increasing the fixed precision pays their constant sum. This is why
\(J\) is chosen after the mesh, conductor, and pattern parameters.

The clean and exceptional parts sum to the literal operator \(E_o\).
Apply their squared estimates at precision \(E+1\). The inequality
\(\|u+v\|^2\le2\|u\|^2+2\|v\|^2\) gives

\[
\mathcal E(E_oF)\le4L^{-2(E+1)}\mathcal E(F)
                  \le L^{-2E}\mathcal E(F). \tag{T18}
\]

The exceptional estimate did not require cancellation. The same positive
row and column bounds, without the rare-event factor, apply if an
arbitrary coefficient of modulus at most one is attached to each
expanded transition. Thus, for a fixed exponent \(\kappa\),

\[
\mathcal E(E_o^\phi F)\le L^{2\kappa}\mathcal E(F),
\qquad |\phi|\le1. \tag{T19}
\]

The exponent is independent of \(J\); its formula is given below. This
distinction is needed in the forest proof: a sign multiplier at a
collision owner uses (T19), not (T18).

\Needspace{6\baselineskip}

\hypertarget{marked-correlation-background-transitions-and-the-constants-in-the-moment-bound}{%
\subsubsection{Background transitions and the constants in the moment
bound}\label{marked-correlation-background-transitions-and-the-constants-in-the-moment-bound}}

At visit \(i\), the background operator \(G_i\) leaves the SMALL state
and all active BIG marks fixed. In one BIG group let \(Z\) be the
indices of OLD atoms at residue zero. Choose \(A\subset Z\), delete
those OLD occurrences, and append \(b\) new atoms
\((p_1,0),\ldots,(p_b,0)\). The coefficient of this choice is exactly

\[
 \left(-\frac{\eta_i}{\rho_0}\right)^{|A|}
 \left(\frac{q_i}{\rho_0^2}\right)^{|Z|-|A|}
 \frac1{b!}\left(-\frac{\eta_i V_g}{\rho_0}\right)^b
 \prod_{j=1}^b\widetilde\mu_g(p_j).
\tag{T20}
\]

Sum over these choices, with the finite output cap, then take the
product over groups and the orbit average (T4). These appended
occurrences are \emph{background births}. A later use of an existing
occurrence is not a birth. Formula (T20) explains both the signs and the
single birth factorial. The OLD factorial in (T3) is the Hilbert measure
and is not inserted a second time in an edge coefficient. The powers of
\(\rho_0\) in (T5) and (T20) cancel at successive endpoints; they are an
auxiliary norm device, not a change to the original residue density.

Here are explicit constants for (F7)--(F11). They also distinguish the
row constant used in removing the OLD cutoff from the smaller
conditional-map constant (T16). For \(d_0=\ell+m\) put

\[
\begin{aligned}
C_{\rm prom}(\rho_0,d_0)&=\sum_{j=0}^{d_0}\frac{j!}{(1-\rho_0)^{j+1}},\\
C_{\rm small}&=\prod_{g\in\mathrm{SMALL}}
       \frac{\ell!}{V_g^\ell(1-\sqrt q)^{\ell+1}},\\
S_{\rm small}&=\sum_{x\in\prod_{g\in\mathrm{SMALL}}\mathscr S_g}
                    \prod_g\frac1{|\mathbb T_g|V_g^M},\\
a&=\mathcal K C_{\rm small}
 \prod_{g\in\mathrm{BIG}}
 \left[w_0^{d_0}4^{d_0}C_{\rm prom}(\rho_0,d_0)
                 \max(1,V_g^{-1})^{d_0}2^{d_0}\right].
\end{aligned}\tag{T21}
\]

Here \(\mathcal K\) is (T11), and \(w_0\) is specified in
\protect\hyperlink{marked-correlation-restoring-all-histories}{the
history restoration}. Thus \(S_{\rm small}\) is a finite total mass, not
a universal constant or a probability normalization. Both
\(C_{\rm small}\) and \(a\) depend on the actual group masses and \(L\);
their polynomial \textbf{exponents} are fixed. The bounds on \(V_g\) and
the number of groups make \(\mathcal K C_{\rm small}\le L^{C_E'}\) and
\(a\le L^{C_b}\) eventually. \(C_E'\) here is a row-budget exponent,
unrelated to the ordinary-slot exponent \(C_E\) in (B1).

An explicit exponent for (T19) can be read without choosing \(J\).
Choose \(C_G\) so that both group counts and
\(|\mathcal A|+|\mathcal B|\) are at most \(C_G\log L\) and put

\[
\begin{aligned}
R_v&=\max(1,v_{\max},v^{-1}),\\
B_{\rm row}&=\max\{1,4^{d_0}C_{\rm prom}(\rho_0,d_0)
                                  \max(1,v^{-1})^{d_0}2^{d_0}\},\\
\kappa&=mC_G\log2+4C_q+2C_\xi+1\\
&\quad+C_G\log\max\left(1,
             \frac{\ell!}{(1-\sqrt q)^{\ell+1}v^\ell}\right)
       +C_G\log B_{\rm row}+\frac{d_0C_G}2\log R_v.
\end{aligned}\tag{T22}
\]

To see the uniformity, at most \(d_0\) slots per group are unshared. If
\(z\) OLD occurrences have residue zero, choosing \(j\le d_0\) of them
costs at most \((z)_j\rho_0^{z-j}\le C_{\rm prom}\). The SMALL row is
bounded by the same falling-factorial calculation as (T16), with
\(\sqrt q\) in place of \(q\). The number of patterns and the character
caps give the first line of the exponent in (T22). The comparison of
forward and reverse harmonic masses contributes at most \(R_v^{d_0G}\)
to the squared bound. Weighted Schur, followed by (T4), gives (T19).
Applying the rare-row estimate (M12) to this same calculation gives
(M13), so \(J\) can indeed be chosen afterwards.

The background estimate is uniform even in unit multipliers on its
individual choices. For \(W_0,w_0,\rho_0\) as below, summing (T20) in
absolute value with weight \(w_0^{|\mathrm{OLD}|}\) uses

\[
\frac1{\rho_0w_0}+\frac{\sqrt q}{\rho_0^2}\le1,
\qquad \sum_{p\in P_g}\widetilde\mu_g(p)\le2.
\]

The deletion sum is binomial and the birth sum is exponential. In the
transposed sum the factorials in (T3) give the corresponding binomial
coefficients; hence the same weighted Schur argument gives
\(\mathcal E(G_i^\phi F)\le\Gamma^2\mathcal E(F)\), with
\(\Gamma=\exp(4W_0w_0G/\rho_0)\). Set

\[
l=L^{-S_0},\qquad h=L^\kappa,\qquad
b=L^{S_0+\kappa},\quad K_r=S_0+\kappa,\qquad
P_0=\left|\prod_{g\in\mathrm{BIG}}\prod_{p\in P_g}b_p\right|
                \mathcal E(v)\Gamma^{N+1},
\tag{T23}
\]

where \(v\) is one on states with empty OLD memory and zero otherwise.
The initial distinct-active mask is retained in the history sum and can
only reduce its boundary energy. Equations (T18)--(T23) supply all three
estimates in (F7); the forest argument below restores the omitted
distinctness between different birth times.

\Needspace{6\baselineskip}

\hypertarget{marked-correlation-restoring-all-histories}{%
\subsection{Restoring all
histories}\label{marked-correlation-restoring-all-histories}}

Equalities between prime labels in a history are indexed by forests. The
exact signed expansion, its coefficients, and the two bounds below are
proved in
\protect\hyperlink{forest-moments}{Section~\ref{forest-moments}}. Its
random-sign factorization explains why each identification loses the
local contraction at at most two visits. For
\(N=2\lceil L^{13/25}\rceil\), the number of available labels is bounded
by

\[
\mathcal A=(N+1)\#\mathrm{BIG}(J+1+\lceil L^2\rceil)\le L^4.
\]

A forest of rank \(r\) requires \(r\) identifications and there are at
most \(\mathcal A^{2r}\) such forests. The expansion has two bounds for
every forest contribution \(F_r\):

\[
|F_r|\le
\min\{P_0(L^{-S_0})^N b^{2r},\ P_0a^N\delta^r\}, \tag{M15}
\]

where

\[
\begin{gathered}
P_0\le L^{(C_v+2)N},\qquad b\le L^{K_r},\qquad a\le L^{C_b},\\
S_0=E+C_v+5,\qquad \delta=\frac2{v e^{L^{3/10}}}.
\end{gathered}
\]

The first bound uses the local signed contraction at each step; each
identification replaces at most two local factors by their row bounds.
The second uses absolute row bounds everywhere and charges one factor
\(\delta\) to each independent equality of BIG prime labels. This is a
finite reindexing of the same history expansion.

Split at \(r=\lceil NL^{-3/20}\rceil\). Below the split, the exponent
cost \((8+2K_r)r\) from the forest count and \(b^{2r}\) is at most \(N\)
eventually. The first bound in (M15) is therefore at most
\(L^{-(E+2)N}\) after counting forests of that rank. Above the split,

\[
(\mathcal A^2\delta)^r
 \le e^{-rL^{3/10}/2}
 \le e^{-NL^{3/20}/2},
\]

which dominates \(L^{(C_v+C_b+E+4)N}\). The second bound gives the same
\(L^{-(E+2)N}\). There are at most \(\mathcal A+1\le L^N\) ranks;
summing proves an upper bound \(L^{-EN}\) for the signed history sum
with the OLD cutoff \(\#\mathrm{OLD}\le\lceil L^2\rceil\).

There is no circular parameter choice here. The fixed volume exponent
\(C_v\) is chosen from the harmonic mass bounds and \(q\), before \(m\)
or \(J\). One may first choose \(\rho_0\in(q^{1/4},1)\) and then use

\[
W_0=\max(1,v_{\max}),\quad
w_0=\max\!\left(1,\frac{4W_0}{\rho_0(1-\sqrt q/\rho_0^2)}\right),
\quad C_v=\frac{4W_0w_0C_0}{\rho_0},
\]

where \(C_0\log L\) bounds the group count. Thus \(S_0\) is fixed before
the local parameters. The later exponents \(K_r,C_b\) affect only the
eventual threshold in the rank split.

To remove this cutoff, put \(T_0=\lceil L^2\rceil\) and
\(H_b=\lceil L^2/4\rceil\). The number of active BIG labels encountered
at all visits is at most \((N+1)G(J+1)\), where \(G\) is the BIG group
count. Eventually \(H_b+(N+1)G(J+1)\le T_0\). The first crossing of the
OLD cutoff therefore requires at least \(H_b\) background births. In the
positive history expansion, multiply each birth weight by \(2\). The
discarded mass is at most \(2^{-H_b}\) times this single positive root
sum; no factor counting histories is needed. Its explicit bound is

\[
S_{\rm small}e^{4(N+1)(1+\sqrt L)}
 \left[C_\nu C_{\rm small}
  \left(16\max(1,v^{-1})e^{\sqrt L}\right)^{(1+m)G}\right]^N,
\]

Here \(C_\nu=\mathcal K\), and

\[
C_\nu C_{\rm small}\le L^{C_E'},\qquad
S_{\rm small}\le
 \exp\!\bigl((C_0+1)L\log L(\sqrt L+\log\max(1,v^{-1}))\bigr).
\] The discarded mass is consequently at most
\(\exp(-\tfrac14(\log2)L^2+O_{\rm fixed}(L^{7/4}))\), hence at most
\(L^{-(E+1)N}\). Use the forest bound at precision \(E+1\) as well.
Adding the two errors gives \(2L^{-(E+1)N}\le L^{-EN}\) for the full
history sum.

\Needspace{6\baselineskip}

\hypertarget{marked-correlation-from-residue-averages-to-the-integer-interval}{%
\subsection{From residue averages to the integer
interval}\label{marked-correlation-from-residue-averages-to-the-integer-interval}}

For a fixed history \(\omega\), let \(h_i\) be its offsets. Its density
at an integer \(n\) is

\[
D_\omega(n)=\prod_p\prod_{i=0}^N
\begin{cases}
\mathbf1_{p\mid n+h_i},&p\text{ is active at visit }i,\\
q_i^{\mathbf1_{p\mid n+h_i}},&\text{otherwise},
\end{cases} \tag{M16}
\]

with \(q_0=q_N=\sqrt q\) and \(q_i=q\) internally. Its model density
\(v_\omega\) is the product of the exact averages over residues modulo
each prime. This retains coincidences between offsets modulo a prime.
Explicitly, let \(A_p\) be the visits at which \(p\) is active, and
\(I_{p,t}=\{i:h_i\equiv t\pmod p\}\). The local factor is

\[
v_{\omega,p}=\begin{cases}
0,&\text{the offsets on }A_p\text{ are not all congruent},\\
\displaystyle\frac1p\prod_{i\in I_{p,t}\setminus A_p}q_i,
 &A_p\ne\varnothing,\ h_i\equiv t\ (i\in A_p),\\
\displaystyle1-\frac1p\sum_{t\bmod p}
              \left(1-\prod_{i\in I_{p,t}}q_i\right),&A_p=\varnothing.
\end{cases}\tag{M16a}
\]

The minus sign between the initial integer residue and \(t\) disappears
under uniform averaging. There is one residue per \textbf{numerical
prime} in (M16a), even if it appears in many active lists.

We verify its connection to the memory expansion, which is needed before
applying an interval estimate. A prime lifespan is either empty or an
interval \([a,z]\subset[0,N]\). For a nonempty lifespan require
\(h_a\equiv h_z\pmod p\), require every active visit to lie in \([a,z]\)
and have offset congruent to \(h_a\), and allow \(a=z\) only when \(a\)
is active. Its weight is

\[
\frac1p\epsilon_a\epsilon_z
 \prod_{a<i<z,\ i\notin A_p}
       (1-\eta_i1_{h_i\equiv h_a\ (p)}),
\qquad
\epsilon_i=\begin{cases}1,&i\in A_p,\\-\eta_i,&i\notin A_p.\end{cases}
\tag{T24}
\]

The empty lifespan is allowed only if \(A_p=\varnothing\) and has weight
\(b_p=1-\sum_i\eta_i/p\). The sum of all these weights is exactly
(M16a). For an absent prime, expand the elementary identity

\[
\prod_{i\in I}(1-\eta_i)
 =1-\sum_{i\in I}\eta_i
  +\sum_{a<z,\ a,z\in I}\eta_a\eta_z
                      \prod_{a<i<z,\ i\in I}(1-\eta_i)
\tag{T25}
\]

in each residue class \(I=I_{p,t}\) and average over \(t\). The identity
follows by telescoping twice, or induction on \(|I|\). For a represented
prime, let \(a_0,z_0\) be its first and last active visits. The inactive
factors strictly before \(a_0\) have the expansion
\(1-\sum_{a<a_0}\eta_a1_{h_a\equiv h_{a_0}}  \prod_{a<i<a_0}(1-\eta_i1_{h_i\equiv h_{a_0}})\);
use the reversed expansion after \(z_0\) and leave the middle factors
unchanged. Multiplication gives exactly (T24), including the cases
\(a=a_0\), \(z=z_0\), and a single active visit.

To recover the transition description from these lifespans, follow each
prime through its active visits and keep it in OLD between separated
active runs. Leaving an active run stores it; returning promotes the
same occurrence. An inactive left endpoint of a lifespan is a background
birth, and an inactive right endpoint a background deletion. The
intermediate residue is \(h_i-h_a\pmod p\), so all zero tests are the
ones in (T24). Conversely, a history with distinct births determines
these lifespans and active runs uniquely, apart from the ordering of
simultaneous births and OLD representatives. Their permutation counts
cancel respectively the birth factorials in (T20) and the sector
factorials in (T3). At each prime the \(\rho_0\) powers cancel;
\(b_p\widetilde\mu_g(p)=1/(V_gp)\); the remaining \(V_g^{-1}\) factors
count active runs. These are exactly the factors obtained from the
integer matrix's state and edge masses. Thus multiplication over primes
and summation over the frame choices identify the restored signed memory
sum with \(\sum_\omega c_\omega v_\omega\). The frame coefficient
retains its sign, auxiliary ban and complex phase throughout this
equality.

Here is the finite approximation used for (M16). If active congruences
at some prime are incompatible, both the integer density and its residue
average vanish. Otherwise let \(Q_0\) be the product of the distinct
active primes. Their congruences specify one class modulo \(Q_0\), with
a fixed positive factor from the inactive visits at those primes. For
each remaining prime put

\[
f_p(r)=1-\prod_{i=0}^N q_i^{1_{r+h_i\equiv0\ (p)}}.
\]

Then \(0\le f_p\le1\), it is supported on at most \(N+1\) residue
classes, and its average is at most \((N+1)/p\le1/2\) eventually, since
every marked prime is at least \(e^{L^{1/10}}\). Put

\[
\lambda=\sum_p\frac1p\sum_{r\bmod p}f_p(r),\qquad
V_0=\prod_p\left(1-\frac1p\sum_r f_p(r)\right),\qquad
U_0=1+\sum_p\sum_r f_p(r).
\]

Bonferroni truncation at orders \(2K\) and \(2K+1\), followed by CRT on
each retained product, gives relative error on an interval of length
\(X\) at most

\[
\frac{e^{4\lambda}}{2^{2K+1}}
 +\frac{Q_0}{X}\,2(2K+2)e\,U_0^{2K+1}e^{2\lambda}. \tag{M17}
\]

Indeed, the omitted elementary symmetric sums are bounded using the sum
\(\lambda\) of the local averages; the normalization by \(V_0\) costs at
most \(e^{2\lambda}\). For each retained residue pattern, CRT counts its
interval representatives with error at most one. Summing the absolute
coefficients of the retained patterns gives the second term in (M17).

Take \(K=\lceil L^{11/20}\rceil\). The actual ranges give

\[
\lambda=O(N\log L),\qquad
\log Q_0=O_J(L^{23/25}\log L),\qquad
K\log U_0=O(L^{19/20}),
\]

whereas \(X\ge e^{\gamma L}\) for a fixed \(\gamma>0\). Both terms in
(M17) are smaller than \(L^{-EN}\) for every fixed \(E\) eventually. For
complex history coefficients, the transfer is made before discarding any
signs:

\[
\left|\sum_\omega c_\omega\sum_nD_\omega(n)
 -X\sum_\omega c_\omega v_\omega\right|
\le XL^{-EN}\sum_\omega|c_\omega|v_\omega. \tag{M18}
\]

The absolute mass on the right is at most \(L^{CN}\). Its exponent is
fixed after \(J\); increasing the requested CRT accuracy absorbs it. The
signed main sum on the left is the one already bounded by (M15).

We spell out the finite physical matrix to which these history estimates
apply. Fix a seed interval \(I=[a,b)\cap\mathbb Z\) of positive
integers, of length \(X=b-a\ge e^{\gamma L}\). Put
\(R=\lceil L^{13/25}\rceil\),
\(r_{\rm path}=\lceil2R\Delta_{\max}\rceil\), with \(\Delta_{\max}\) as
in (M20a), and take the carrier positions to be

\[
I^+=[a-r_{\rm path},b+r_{\rm path}]\cap\mathbb Z.
\]

Every prefix of every \(2R\)-step history from \(I\) belongs to \(I^+\):
each displacement has absolute value at most \(\Delta_{\max}\), so the
triangle inequality bounds a prefix by \(2R\Delta_{\max}\). Thus
truncating the matrix to this common carrier introduces no
history-dependent restriction on its starting interval. All residue
counts below use the same \(I\) and the same length \(X\).

Put \(M=J+1\). A state is an integer position \(n\in I^+\), together
with, for each marked group \(g\), an ordered list \(\boldsymbol p_g\)
of \(M\) distinct primes dividing \(n\). The residues of \(n\) determine
all residue coordinates of the state. All states have the same reference
mass

\[
m_0=\prod_g V_g^{-M};\qquad
u_n=\sqrt{m_0}\sum_{\substack{s\text{ a state}\\\operatorname{pos}(s)=n}}e_s.
\tag{M18a}
\]

Thus no mark fibre is normalized to have mass one.

A frame consists of a raw pattern or a nonempty pair \((a,b)\) from
(M2), an auxiliary prime for each freed BIG probe coordinate, and
independent source and target permutations \(\sigma_g,\tau_g\in S_M\) in
every group. For a comparison pattern, let
\(C_g\subseteq\{0,\ldots,m-1\}\) be its freed coordinates in group
\(g\); for a raw pattern and for SMALL groups, \(C_g=\varnothing\). The
shared coordinates are \(S_g=\{0,\ldots,J-1\}\setminus C_g\). Write the
auxiliary primes as \(\pi_{g,i}\) and give them the law

\[
\mu_{\rm aux}=\prod_{g\in\mathrm{BIG}}\prod_{i\in C_g}
                 \frac1{V_g\pi_{g,i}}.
\]

Auxiliary values may repeat. Each must differ from every source and
target mark in its own group. For states \(s,t\), a frame is allowed
when their permuted marks agree at each shared coordinate and

\[
\begin{aligned}
D_f&=\prod_g\prod_{i\in S_g}p^s_{g,\sigma_g(i)}
       \prod_{g\in\mathrm{BIG}}\prod_{i\in C_g}\pi_{g,i},\\
\operatorname{pos}(t)&=\operatorname{pos}(s)+2kD_f.
\end{aligned} \tag{M18b}
\]

All factors of \(D_f\) retain their multiplicities. For the raw pattern
let \(K_f=1\). For a comparison pattern put

\[
K_f=-(-1)^{|a|+|b|}K_A(D_a,X_a)K_B(D_b,X_b),
\]

where \(D_a,D_b\) are the corresponding products of auxiliary primes,
\(X_a\) is the product of the \textbf{target} marks in the first probe
block, and \(X_b\) is the product of the \textbf{source} marks in the
second block. Both endpoint products use the permuted lists. The two
kernels are exactly (M4), with the actual cell masses for these probe
products.

For an allowed frame define

\[
w_f(s,t)=\mu_{\rm aux}\prod_g
 \left[(M!)^{-2}V_g^{-(M-|S_g|)}
 q^{(\omega_g(\operatorname{pos}(s))+
       \omega_g(\operatorname{pos}(t))-2M)/2}\right].
\]

The powers of \(q\) are the two endpoint dampings. Their exponents are
nonnegative because each endpoint has \(M\) distinct marked divisors.
The physical matrix is

\[
A_{s,t}=\sum_{f\text{ allowed}}w_f(s,t)K_fD_f^{i\zeta}. \tag{M18c}
\]

Reversing the two endpoint permutations preserves \(D_f\) and the real
weight. Reversing the orientation conjugates the whole complex
coefficient, including \(K_f\). The backwards matrix is therefore
exactly \(A^*\).

Expanding \((AA^*)^R\) gives the alternating histories above. Insert
\(\mathrm e(2kD_f\theta)\) on a forward edge and its conjugate on a
backward edge. The product is \(\mathrm e(\theta h_N)\), so integration
for \(0\le\theta\le1\) keeps precisely the histories returning to the
initial integer position. The initial and final mark lists remain free;
the vectors \(u_n\) in (M18a) perform those sums. The restored history
bound and (M18) consequently give

\[
T:=\sum_{n\in I}\langle u_n,(AA^*)^Ru_n\rangle
\le XL^{-2RE},\qquad R=\lceil L^{13/25}\rceil. \tag{M19}
\]

\Needspace{6\baselineskip}

\hypertarget{marked-correlation-the-endpoint-bound}{%
\subsection{The endpoint
bound}\label{marked-correlation-the-endpoint-bound}}

The first endpoint has the form \(f=\sum_{n\in I}c_nu_n\); its
coefficient depends on the integer position but not on the individual
marks. The second endpoint \(g\) may depend on all marks. The finite
spectral estimate used here is

\[
|\langle f,Ag\rangle|
\le\sqrt3\,\|g\|\,\|f\|^{1-1/R}
 (H_{\rm loc}C_{\rm sup}^2)^{1/(2R)}T^{1/(2R)}. \tag{M20}
\]

Here \(|c_n|\le C_{\rm sup}\) and a nonzero matrix entry changes the
integer position by at most the positive integer \(H_{\rm loc}\). To
prove the estimate, partition the positions into disjoint intervals
\(I_j\) of \(H_{\rm loc}\) integers. Let \(f_j\) be the restriction of
\(f\) to \(I_j\), and let \(g_j\) be the restriction of \(g\) to
\(I_{j-1}\cup I_j\cup I_{j+1}\). Locality gives
\(\langle f,Ag\rangle=\sum_j\langle f_j,Ag_j\rangle\) and
\(\sum_j\|g_j\|^2\le3\|g\|^2\). Put \(P=AA^*\). Spectral interpolation
for the nonnegative matrix \(P\), followed by Cauchy--Schwarz, gives

\[
|\langle f_j,Ag_j\rangle|
 \le\|g_j\|\,\|f_j\|^{1-1/R}
            \langle f_j,P^Rf_j\rangle^{1/(2R)}.
\]

Indeed, the spectral measure of \(P\) at \(f_j\) satisfies
\(\langle f_j,Pf_j\rangle\le \|f_j\|^{2-2/R}\langle f_j,P^Rf_j\rangle^{1/R}\).
Writing \(f_j=\sum_{n\in I\cap I_j}c_nu_n\) and applying Cauchy--Schwarz
to the vectors \(P^{R/2}u_n\) yields

\[
\langle f_j,P^Rf_j\rangle
 \le H_{\rm loc}C_{\rm sup}^2
       \sum_{n\in I\cap I_j}\langle u_n,P^Ru_n\rangle.
\]

Hölder's inequality in \(j\), with exponents \(2\), \(2R/(R-1)\) and
\(2R\), proves (M20), since \(\sum_j\|f_j\|^2=\|f\|^2\). For \(R=1\),
omit the middle factor and use Cauchy--Schwarz. In the application
\(R>1\).

In the application we may take \(C_{\rm norm}=9\) and
\(C_{\rm sup}=e^{5\sqrt L}\), with \(\|f\|,\|g\|\le\sqrt X L^9\)
eventually. To specify the window, put

\[
\begin{aligned}
C_D&=4(J+1)(4+1/\log2),\\
\Delta_{\max}&=2L^{C_k}\exp(C_D\log L\,L^{2/5}),\\
H_{\rm loc}&=2\lceil2R\Delta_{\max}\rceil+1.
\end{aligned}\tag{M20a}
\]

Every frame satisfies \(D_f\le\exp(C_D\log L\,L^{2/5})\), including
repeated auxiliaries, and every \(2R\)-step path from \(n\) stays within
\(\lceil2R\Delta_{\max}\rceil\) of \(n\). The reference measure in (M20)
is ordinary counting measure on the finite physical states;
\(\sqrt{m_0}\) is in the vectors, not the norm. Since
\(\log H_{\rm loc}=O_J(\log L\,L^{2/5})\) and \(R\ge L^{13/25}\), the
local factor in (M20) is bounded. Apply (M19) with
\(E'=E+2C_{\rm norm}+1\). Then (M20) is at most \(XL^{-E}\) eventually.
The \protect\hyperlink{physical-lift}{physical-lift argument} gives the
exact endpoint vectors, the Fourier pairing identity (L8), and the
shared-prime restoration error (L4). It obtains the \(J+1\) marked
physical states by multiplying both original integers by a common
product of \(J\) primes per group; it does not require the original
integers in (B17) to have that many distinct marked divisors.

This proves the saving for the original term minus its arithmetic
comparison term. It remains to estimate that comparison term itself.

\Needspace{6\baselineskip}

\hypertarget{marked-correlation-the-comparison-term-and-the-fine-discrepancy}{%
\subsection{The comparison term and the fine
discrepancy}\label{marked-correlation-the-comparison-term-and-the-fine-discrepancy}}

We give the endpoint used in the comparison explicitly. If \(G\) is the
union of the marked prime groups, let \(n_*\) be the part of \(n\) with
all prime factors in \(G\) removed. For bounded coefficients
\(\alpha,c\), a prime slot \(Q_*\) disjoint from \(G\), a character
\(\chi\), and real frequencies \(\sigma_0,\sigma_1\), set

\[
\operatorname{Core}(n_*)=
 \sum_{\substack{m p\mid n_*\\p\in Q_*}}
 \alpha(m)\chi(m)m^{i\sigma_0}p^{i\sigma_1}c(n_*/mp).
 \tag{M21}
\]

For each marked group \(P_g\), let \(\mathcal L_g(n)\) consist of its
ordered lists of \(t_g\) distinct prime divisors of \(n\), and let
\(\omega_g(n)\) count those distinct divisors. For a function \(b\) of
the BIG lists with \(|b|\le1\), the literal marked weight is

\[
W_{\boldsymbol t}^{,b}(n)=
 \prod_g\frac{q^{\omega_g(n)-t_g}}{V_g^{t_g}}
 \sum_{\substack{\boldsymbol p_g\in\mathcal L_g(n)\\\text{all }g}}
 b((\boldsymbol p_g)_{g\in\mathrm{BIG}}). \tag{M22}
\]

This includes the distinguished SMALL group; in the source its number of
marks is \(k+1\). Multiplicities in \(n\) are not restricted by the
distinctness of the marks. With a smooth cutoff \(\Lambda(n)\), the
endpoint Fourier polynomial is

\[
\widehat E(\theta)=\sum_{1\le n\le N_{\rm cap}}
 \Lambda(n)\varepsilon(n)\operatorname{Core}(n_*)
 W_{\boldsymbol t}^{,b}(n)\mathrm e(n\theta). \tag{M23}
\]

The cell/character kernel (M4) has a concrete rank decomposition: each
\((j,\chi)\) contributes the product of
\(1_{D\in I_j}\overline{\chi(D)}\) and \(1_{X\in I_j}\chi(X)\), with
coefficient \(\nu_j^{-1}\). There are at most \(n_AQ^2\) such terms in
the first kernel and \(n_BQ^2\) in the second. Thus a comparison pattern
has at most \(n_An_BQ^4\) pairs, each coefficient at most \(\xi^{-2}\).
The support counts (M9a) bound this rank by a fixed power of \(L\). For
retained cells there is a sharper bound: their disjoint probabilities
sum to at most one, so their number is at most \(\xi^{-1}\) in either
block. Including all patterns gives

\[
\begin{aligned}
\#\{\text{ranks}\}&\le L^{C_{\rm pat}}Q^4\xi^{-2},\\
\sum_{\text{ranks}}\nu_{j_A}^{-1}\nu_{j_B}^{-1}
 &\le L^{C_{\rm pat}}Q^4\xi^{-4},\qquad
C_{\rm pat}=m(3+1/\log2)\log2.
\end{aligned}\tag{M23a}
\]

Here \(3+1/\log2\) is an upper group-count constant for the actual BIG
family. The two inverse cell masses stay outside the endpoint tests. The
remaining endpoint cell indicators and characters are the bounded list
coefficient \(b\) in (M22).

The auxiliary sum retains its dependence on the shift and cutoff. More
explicitly, let \(\nu_A(u),\nu_B(v)\) be the probability laws of the two
auxiliary products. Independence of the two disjoint probe blocks gives
the exact finite identity

\[
\begin{aligned}
&\sum_{\rm aux}\mu_{\rm aux}
 1_{D_A\in I_{j_A},D_B\in I_{j_B}}
 \overline{\chi_A(D_A)\chi_B(D_B)}\,H(D_AD_B)\\
&\quad=\sum_{u\in I_{j_A}}\sum_{v\in I_{j_B}}
 \nu_A(u)\nu_B(v)\overline{\chi_A(u)\chi_B(v)}\,H(uv).
\end{aligned}\tag{M23b}
\]

The inverse cell masses and the pattern sign multiply both sides. After
the shared-divisor restoration described below, \(H(c)\) retains the
shift \(2kc\) and the cutoff \(\Psi(n/(Zc))\). Thus (M23b) does not
replace this dependence by a mass-one average. Fourier inversion of the
cutoff as in (L8) separates its Mellin factor; additive Fourier
inversion separates the two endpoint integers. This yields the pair of
polynomials (M23), with an auxiliary bilinear phase
\(\mathrm e(-2kuv\theta)\) on the right of (M23b). Partition \(u,v\)
into the coarse dyads of (M9a) when applying (M7). The scales of a
retained cell pair are \(H'=e^{j_Ah_A+j_Bh_B}\) and \(Y=ZH'\); its
actual product lies between \(H'\) and \(e^{h_A+h_B}H'\). A fixed smooth
cutoff equal to one on this ratio interval gives the endpoint support
used in (M25). All rank, coefficient and Fourier costs are fixed before
requesting the discrepancy precision below.

Here is the resulting comparison identity with the endpoint factors
displayed. For a pattern \((a,b)\) let \(t_g=M-|S_g|\) be its number of
free marks; thus \(t_g=1\) in SMALL and \(t_g=1+|C_g|\) in BIG. For rank
\(r=(j_A,\chi_A,j_B,\chi_B)\) define the two list functions

\[
b_{1,r}=\overline{1_{X_B\in I_{j_B}}\chi_B(X_B)},\qquad
b_{2,r}=1_{X_A\in I_{j_A}}\chi_A(X_A).
\tag{M23c}
\]

Here \(X_B\) is the source probe product and \(X_A\) the target probe
product, with the shared entries deleted and the free entries kept in
their original coordinate order. Both functions have modulus at most one
and contain no auxiliary variables. With the two cores of (L7), define

\[
U_r(n)=\varepsilon(n)n^{it_1}\operatorname{Core}_1(n_*)
                  W_{\boldsymbol t}^{,b_{1,r}}(n),\qquad
V_r(n)=\varepsilon(n)n^{it_2}\operatorname{Core}_2(n_*)
                  W_{\boldsymbol t}^{,b_{2,r}}(n).
\tag{M23d}
\]

The first core is conjugated when \(U_r\) is conjugated in the pairing.
For (B17) the phases are \(t_1=-(\tau+\xi_3)\), \(t_2=-\tau\),
\(\zeta=-\xi_3\), as in (L7)--(L8). Write \(c=uv\) and
\(s_{a,b}=-(-1)^{|a|+|b|}\). The independent-rank main contribution of
this pattern and rank to the signed comparison part of \(A\) (hence to
\(-Q_{\rm cmp}\) in (L8)) is exactly

\[
\begin{aligned}
\frac{s_{a,b}}{\nu_{j_A}\nu_{j_B}}
\sum_{\substack{u\in I_{j_A}\\v\in I_{j_B}}}
 &\nu_A(u)\nu_B(v)\overline{\chi_A(u)\chi_B(v)}\,c^{i\zeta}\\
&\cdot\sum_{\substack{n>0\\n+2kc>0}}
 \frac{\Psi(n/(Zc))}{n}\overline{U_r(n)}V_r(n+2kc).
\end{aligned}\tag{M23e}
\]

To check the removal of the shared divisor \(d\), substitute the
original positions \(dn,dn'\) and \(D_f=dc\). The displacement becomes
\(n'=n+2kc\) and the cutoff becomes \(\Psi(n/(Zc))\). The ordinary parts
of the two integers are unchanged, since all prime factors of \(d\) are
marked. The common parity is squared. The phase in \(d\) is
\(d^{i(-t_1+t_2+\zeta)}=1\). The shared harmonic factors then sum to one
on the unrestricted shared family; the exclusions and repeated-mark
cases produce the weighted restoration error (M28), not an exact
mass-one assertion on the restricted family.

There is one further explicit difference in (M23e): the original
auxiliary ban forbids a free endpoint mark from equalling an auxiliary
prime. Partitioning the complete pair of free mark lists into the
allowed and forbidden sets gives, exactly, the independent product of
the two weights minus its collision contribution. A forbidden pair
forces \(p\mid n\) for an auxiliary prime \(p\ge e^{L^{3/10}}\); if it
divides the other endpoint, it also divides \(n\), since \(p\mid c\).
Reindex \(n=pr\). The factor \(1/n\) and the divisor-moment bound for
the two cores give a factor \(1/p\) with a fixed logarithmic cost. If
\(L^{E_{\rm mark}}\) bounds either full marked weight, the error for one
pattern and one fixed auxiliary assignment is at most

\[
C_JK_\Psi Q^4\xi^{-2}
 L^{2E_{\rm mark}+16+C_\psi}e^{-L^{3/10}}.
\tag{M23f}
\]

Here \(C_J=2J(3+1/\log2)\,64\cdot66\cdot3^{15}\); \(E_{\rm mark}\)
depends only on \(J\). Summing auxiliary probabilities and at most
\(L^{C_{\rm pat}}\) patterns preserves an error smaller than every fixed
inverse power of \(L\). Thus (M23e), summed over ranks and patterns,
equals that signed comparison up to the two specified errors (M28) and
(M23f).

For completeness, the Fourier pairing of (M23e) also has an exact form.
Put \(\psi_{\log}(t)=\Psi(e^t)\). Choose a fixed real smooth cutoff
\(\eta\), equal to one at both \(n/Y\) and \((n+2kc)/Y\) whenever the
summand in (M23e) is nonzero, and supported in a fixed compact
subinterval of \((0,\infty)\). Such a choice is uniform because
\(1\le c/H'\le e^{h_A+h_B}\) and \(|k|/Z\to0\) uniformly. Set

\[
A_{r,s,Y}(n)=\frac{\eta(n/Y)}{n}n^{-2\pi is}U_r(n),\qquad
B_{r,Y}(n)=\eta(n/Y)V_r(n),
\tag{M23g}
\]

and use \(\widehat A(\theta)=\sum_n A(n)\mathrm e(n\theta)\), likewise
for \(B\). Then (M23e) is

\[
\int_{\mathbb R}\widehat\psi_{\log}(s)
 \int_0^1\mathcal K_{r,s}(\theta)
       \overline{\widehat A_{r,s,Y}(\theta)}
       \widehat B_{r,Y}(\theta)\,d\theta\,ds,
\tag{M23h}
\]

where the full coupled kernel is

\[
\mathcal K_{r,s}(\theta)=
\frac{s_{a,b}Z^{-2\pi is}}{\nu_{j_A}\nu_{j_B}}
\sum_{\substack{u\in I_{j_A}\\v\in I_{j_B}}}
\nu_A(u)\nu_B(v)\overline{\chi_A(u)\chi_B(v)}
(uv)^{i\zeta-2\pi is}\mathrm e(-2kuv\theta).
\tag{M23i}
\]

The circle integral imposes \(n'=n+2kuv\), and the Mellin integral
restores \(\Psi(n/(Zuv))\), proving (M23h) term by term. The cell masses
in the denominator cancel the total restricted product probability, so
\(|\mathcal K_{r,s}(\theta)|\le1\) for every \(s,\theta\). On minor arcs
(M7)--(M9) give the stronger saving, retaining both inverse cell masses
in the coefficient budget (M23a). The endpoints (M23g) are now
independent of \(u,v\), which is the precise interface needed for the
arc-energy estimate.

The coefficient \(\alpha\) has the complete-cell discrepancy (C3): for
any prescribed \(B\), the fine mesh can be chosen so that

\[
\left|\sum_{a\le m\le b}
       \frac{\alpha(m)\chi(m)m^{it}}m\right|\le L^{-B},
\qquad q,|t|\le L^B,\quad [a,b]\subset[1,x^2]. \tag{M24}
\]

Here \(\alpha\) already includes its restriction to the factor interval
\(K\subset[x^\beta,x^2]\) of multiplicative width at most two, with
fixed \(\beta>0\). Thus the displayed sum means its intersection with
\(K\). If nonempty, that intersection satisfies the lower endpoint and
width conditions of (C3). The modulus \(q\) in this display is an
integer conductor variable, distinct from the fixed damping \(q=1/2\).

Let the two nonempty probe-product scales be \(U_1,U_2\), put
\(H'=U_1U_2\) and \(Y=ZH'\), and suppose \(Z\ge xL^{-C_Z}\). The marked
arc estimate obtained from (M24) is

\[
\int_{|\beta|\le L^{C_{\rm arc}}/H'}
 |\widehat E(h/d+\beta)|^2\,d\beta
 \le K_{\rm arc}YL^{-T}. \tag{M25}
\]

Here \(Q_*\) is the complete prime interval used by the comparison,
contained in \([x^\tau,x^\eta]\) with \(0<\tau<\eta<1/4\); the support
of \(\alpha\) is \(e^{\sqrt L}\)-rough. The cutoff has the prescribed
polynomial bounds on its logarithmic derivatives, and
\(2\le H'/L^{C_{\rm arc}}\le Y/(8A_\psi)\). These conditions hold for
the actual data from (B16). Under them, (M25) is uniform for the
polynomial conductor, frequency, and coefficient budgets in the
comparison expansion. The support has \(e^{3L/4}\le Y\le e^{3L/2}\);
take \(N_{\rm cap}=\lceil A_\psi Y\rceil\le x^2\), where \(A_\psi\)
bounds the cutoff support.

Here are the estimates establishing (M25). Peeling one distinguished
SMALL mark leaves the core and parity at the original integer and
evaluates only the remaining marked weight at \(n/p\). Separate the
smooth logarithmic factors by Fourier inversion, and write
\(H=H'/L^{C_{\rm arc}}\). The finite local-energy inequality (E1) and
the exact separation (E9) in
\protect\hyperlink{local-energy}{Section~\ref{local-energy}} give the
following four terms. With \(B_{\rm amp}\ll L^{E_{\rm amp}}\) and
\(B_s\ll L^{E_s}\) the resulting amplitude bounds, those terms, up to
fixed constants, are

\[
\begin{aligned}
&YV_s^{-2}d^2(1+\log_2N_{\rm cap})^8
                   L^{-(T+2c_m+8)},\\
&Y^{-1}L^{-(j-2)}
 \bigl(LY/H+16N_{\rm cap}(1+\log N_{\rm cap})\bigr)
 B_{\rm amp}^2N_{\rm cap}(1+\log N_{\rm cap})^{63},\\
&(H/Y)^2 B_{\rm amp}^2N_{\rm cap}(1+\log N_{\rm cap})^{63},\\
&3L B_s^2N_{\rm cap}(1+\log N_{\rm cap})^{16}e^{-L^{a_s}}.
\end{aligned} \tag{M26}
\]

The first is the finite Mellin-window contribution (E13), obtained from
(M24) and the small- and long-prime estimates, with \(d\le L^{c_m}\).
The second is the weighted Mellin tail; choose
\(j=\lceil2E_{\rm amp}+64+T\rceil+2\). The third is the error (E3) in
replacing linear windows by logarithmic windows. Since \(Y=ZH'\),
\(H/Y=1/(ZL^{C_{\rm arc}})\), so its exponential decay pays for every
displayed logarithmic factor. The fourth accounts for squares of a
peeled SMALL prime, whose size is at least \(e^{L^{a_s}}\); here
\(a_s=1/10\) is available. Divisor-moment estimates give the displayed
powers \(8,63,16\). Each line of (M26) is \(O(YL^{-T})\) with the stated
choices, proving (M25). The window replacement follows from the Taylor
estimate (E2) and finite overlap, with no prime distribution estimate on
an interval of length \(H\).

The first Fourier polynomial in the actual comparison has an additional
factor \(1/n\). With its cutoff written as \(\eta(n/Y)\),

\[
\frac{\eta(n/Y)}n=Y^{-1}\frac{\eta(n/Y)}{n/Y}.
\]

The cutoff is supported away from zero. Applying (M25) to \(\eta(u)/u\)
therefore gives energy \(K_{\rm arc}Y^{-1}L^{-T}\) for this first
polynomial; the second has energy \(K_{\rm arc}YL^{-T}\).
Cauchy--Schwarz cancels the two powers of \(Y\). Summing the rational
arcs of the comparison gives \(2K_{\rm arc}L^{2c_m-T}\). Off those arcs,
the auxiliary-product bilinear exponential-sum estimate gives
\(L^{-A_1}\) against the global polynomial endpoint energy. Choose
\(T,A_1\) to pay for the finite rank and pattern budgets. This makes the
whole comparison term smaller than any requested logarithmic power.

Subtracting this separately bounded term from the signed-operator
estimate (M20) bounds the original correlation.

\Needspace{6\baselineskip}

\hypertarget{marked-correlation-caps-shared-divisors-and-the-parameter-order}{%
\subsection{Caps, shared divisors, and the parameter
order}\label{marked-correlation-caps-shared-divisors-and-the-parameter-order}}

The finite comparison first imposes a cap on the marked part and allows
marks shared between the two endpoints. Both modifications are removed
quantitatively. The Rankin estimate from the weight chapter gives

\[
\sum_{\substack{u\le Y\\u_P>\lfloor e^{L/88}\rfloor}}W_1(u)
 \le YL^{(3+1/\log2)\log(6e^{12})}e^{-L^{3/5}/88}. \tag{M27}
\]

Stripping shared divisors has error

\[
O\!\left(L^{8(J+1)\log2+2E_{\rm mark}+17+C_\psi}
                e^{-L^{1/10}}\right). \tag{M28}
\]

The quotient cap is transported through the exact finite identity before
using (M27); its final specialization has shared divisor one and zero
extra marks. Thus no assertion that the marked part is squarefree, or
that a cap is unchanged by arbitrary division, is needed. Both (M27) and
(M28) are negligible at every fixed logarithmic accuracy.

The choices can now be made in order. Fix the desired saving and all
external frequency, shift, and cutoff bounds. Choose the local
contraction precision and \(m\) using (M1)--(M3), then the arithmetic
mesh, conductor, and retained-mass parameters using (M4)--(M10). Choose
\(J\) using (M13). The forest row bounds, total path mass, comparison
rank, Fourier-tail orders, and CRT precision are chosen after these
fixed parameters. Finally choose the discrepancy order \(B\) and the
fine-cell exponent \(S=2(B+1)+4\).

These choices precede the positive lower exponent \(\beta\) of the fine
interval and its positive roughness exponent \(\gamma\). Only the
eventual threshold in \(L\) depends on \(\beta,\gamma\). Combining the
endpoint bound (M20), the comparison estimate (M25), and the two
restoration errors (M27)--(M28) gives, with precisely the coefficients
of (B16)--(B17),

\[
|\mathcal C_{k,Z}|\le ZL^{-D_1}
\]

uniformly for \(0<|k|\le L^{C_k}\), \(xL^{-C_Z}\le Z\le xL^{C_Z}\), the
stated polynomial frequency bounds, and either multiplicity-parity
choice. This is (B18).
